\begin{figure}[htbp]
\centering
\begin{tikzpicture}[
  scale=0.55, transform shape, font=\small,
  %%%%%%%%% struct -> clase UML
  struct/.style={draw, fill=white, rectangle split, rectangle split parts=2,
                 rectangle split part align=center, align=left, inner sep=4pt},
  structm/.style={draw, fill=white, rectangle split, rectangle split parts=3,
                  rectangle split part align=center, align=left, inner sep=4pt},
  %%%%%%%%% trait -> interface UML
  trait/.style={draw, dashed, fill=black!4, rectangle split, rectangle split parts=2,
                rectangle split part align=center, align=left, inner sep=4pt},
  %%%%%%%%% crate -> paquete UML
  tab/.style={draw, minimum width=14mm, minimum height=4mm, inner sep=0pt},
  etiqcrate/.style={anchor=north west, xshift=2.5mm, yshift=-2.5mm,
                    font=\footnotesize\ttfamily},
  %%%%%%%%% relacións
  depende/.style={dashed, -{Latex[open, length=2.4mm, width=2mm]}},
  realiza/.style={dashed, -{Triangle[open, length=3mm, width=3mm]}},
  compon/.style={{Diamond[length=3mm, width=2.2mm]}-},
  etiq/.style={font=\scriptsize, inner sep=1.5pt, fill=white},
  membro/.style={font=\footnotesize\ttfamily}
]

%%%%%%%%% quen recolle a entrada de todas as fontes
\node[structm, minimum width=62mm] (mgr) at (0,0) {
  \nodepart{one}  \texttt{InputManager}
  \nodepart[membro]{two}
        + key\_triggers: InputKeyTriggers\\
        + mouse\_triggers: Option<InputMouseTriggers>\\
        + gpio\_triggers: Option<InputGpioTriggers>\\
        + current / previous: InputState
  \nodepart[membro]{three}
        + update(rl, dt)\\
        + is\_held\_*() / is\_pressed\_*()
};

%%%%%%%%% o estado resultante
\node[structm, minimum width=52mm] (state) at (100mm,0) {
  \nodepart{one}  \texttt{InputState}
  \nodepart[membro]{two}
        + up / down / left / right: bool\\
        + a / b / start / select: bool\\
        + escape / speed\_up: bool
  \nodepart[membro]{three}
        + merge(rl, states)\\
        + apply(joypad)
};

%%%%%%%%% a interface común a todas as fontes de entrada
\node[trait, minimum width=52mm] (tois) at (100mm,-52mm) {
  \nodepart{one}  \scriptsize$\ll$trait$\gg$\\[1pt] \texttt{ToInputState}
  \nodepart[membro]{two}
        + to\_input(rl) -> InputState
};

%%%%%%%%% a fonte de entrada da Gamepi13
\node[structm, minimum width=62mm] (gpiot) at (0,-52mm) {
  \nodepart{one}  \texttt{InputGpioTriggers}
  \nodepart[membro]{two}
        - pins: Arc<GpioPins>
  \nodepart[membro]{three}
        + new()\\
        + shared()\\
        + to\_input(rl)
};

%%%%%%%%% os pines do conector, co seu número BCM
\node[struct, minimum width=62mm] (pins) at (0,-112mm) {
  \nodepart{one}  \texttt{GpioPins}
  \nodepart[membro]{two}
        - up: InputPin (5) \quad - down: InputPin (6)\\
        - left: InputPin (16) \quad - right: InputPin (13)\\
        - a: InputPin (21) \quad - b: InputPin (20)\\
        - x: InputPin (15) \quad - y: InputPin (12)\\
        - l: InputPin (23) \quad - r: InputPin (14)\\
        - start: InputPin (26) \quad - select: InputPin (19)
};

%%%%%%%%% a dependencia externa
\node[struct, minimum width=42mm] (pin) at (100mm,-112mm) {
  \nodepart{one}  \texttt{InputPin}
  \nodepart[membro]{two}
        + into\_input\_pullup()\\
        + is\_low()
};

\begin{scope}[on background layer]
  \node[draw, fill=black!2, fit=(pin), inner xsep=5mm, inner ysep=7mm] (rpipal) {};
  \node[tab, fill=black!2, anchor=south west] at (rpipal.north west) {};
\end{scope}
\node[etiqcrate] at (rpipal.north west) {rpi\_pal};

%%%%%%%%% relacións
\draw[compon] (mgr.east) -- node[etiq, above]{current} (state.west);
\draw[compon] (mgr.south) -- node[etiq, right]{gpio\_triggers} (gpiot.north);
\draw[realiza] (gpiot.east) -- node[etiq, above]{implementa} (tois.west);
\draw[depende] (tois.north) -- node[etiq, right]{$\ll$devolve$\gg$} (state.south);
\draw[compon] (gpiot.south) -- node[etiq, right]{pins} (pins.north);
\draw[depende] (pins.east) -- node[etiq, above, align=center]{12 pines\\ en modo \textit{pull-up}} (rpipal.west);

\end{tikzpicture}
\caption{Recollida da entrada dos botóns da Gamepi13 a través dos \acrshort{gpio}.}
\label{fig:gpio-input}
\end{figure}
